\documentclass[10pt,a4paper]{amsart}
\usepackage[margin=19mm]{geometry}
\usepackage{amsmath,amssymb,mathtools}
\usepackage[scr=rsfs,cal=euler]{mathalfa}
\usepackage{xfrac}
\usepackage[unicode,hidelinks,bookmarks=false]{hyperref}
\newcommand{\ie}{i.e.\ }
\newcommand{\cf}{cf.\ }
\newcommand{\resp}{respectively\ }
\newcommand{\Subsection}{\subsection}
\providecommand{\nobreakdash}{\nobreak\hbox{-}}
\setlength{\emergencystretch}{2em}
\multlinegap0pt
\usepackage{bm}
\usepackage{bbm}
\usepackage{dsfont}
\usepackage{xspace}
\usepackage{cancel}
\usepackage{mathtools}
\usepackage{amsthm}
\makeatletter
\@ifundefined{theorem}{\newtheorem{theorem}{Theorem}}{}
\@ifundefined{claim}{\newtheorem{claim}[theorem]{Claim}}{}
\@ifundefined{lemma}{\newtheorem{lemma}[theorem]{Lemma}}{}
\makeatother
\def\q{\hfill\rule{1ex}{1ex}}
\newcommand{\ber}{\text{Berge-}}
\newcommand{\mg}{\mathcal{G}}
\newcommand{\mh}{\mathcal{H}}
\title[Extremal results on Berge disjoint paths]{Extremal results on Berge disjoint paths}
\author{Xiamiao Zhao, Yiyan Zhan, Mei Lu}
\date{Source: arXiv:2512.23382v2 (CC BY 4.0)}
\begin{document}
\maketitle

\noindent\emph{Source and licence.} Extremal results on Berge disjoint paths. Version arXiv:2512.23382v2 (CC BY 4.0), distributed under the Creative Commons Attribution 4.0 International licence, as recorded in the arXiv licence metadata for this exact version and shown on the abstract page https://arxiv.org/abs/2512.23382v2. Full rights notice: licenses/arxiv-2512.23382-CC-BY-4.0.txt.\par\smallskip

\noindent\textit{Editorial scope.} Counted unit: the Lemma whose source label is \texttt{lem: enough Berge common} in the article (source0.tex lines 831--834, source label \texttt{lem: enough Berge common}), delivered together with the article's own complete original proof of it (source0.tex lines 835--952, one \texttt{proof} environment of 118 lines). No mathematics is written, repaired or re-derived: statement and proof are the article's own text line for line. Required context retained uncounted: lem: the ineuqlaity in sec 3 (2) (lines 355--358); 3-2 (lines 820--823). Every definition, display and label of those lines is retained unchanged. Omitted from this package and not redistributed: 1--354, 359--819, 824--830, 953--1666. Nothing inside a retained excerpt is omitted or altered: every retained body is delivered exactly as the article's lines are written, so no omission receipt of any kind accompanies this package. Citations: the retained excerpts carry no citation command at all, so no bibliography entry of the article is transcribed and no reference is redistributed. Reference adaptation: none was needed and none was made; every reference inside the delivered unit resolves inside it, so no unresolved reference or citation is produced. Printed slips kept literal: no printed word, symbol, hypothesis or number of the counted statement or of its proof was altered. Layout and class: the article's own class is \texttt{article} with packages including color, amsmath, extarrows, geometry, authblk, hyperref, ntheorem, enumitem, graphicx, subcaption, euscript, inputenc, enumerate, amssymb; the wrapper is the lane's pinned \texttt{amsart} editorial wrapper at 10pt on A4, which restates the theorem-like environments the retained text needs and adds the article's own macro definitions that the retained text needs (\texttt{\textbackslash ber}, \texttt{\textbackslash mg}, \texttt{\textbackslash mh}, \texttt{\textbackslash q}), with extra wrapper packages (bm, bbm, dsfont, xspace, cancel, mathtools). The delivered numbering is the unit's own. Assets: no retained span contains a figure environment, a graphics-inclusion command, a PDF-page inclusion, a table or a TikZ picture; no image file is copied into this package, the package declares no asset dependency and no image file is redistributed with it. Licence: version arXiv:2512.23382v2 of this article is distributed under the Creative Commons Attribution 4.0 International licence, as recorded in the arXiv licence metadata for this exact version and shown on the abstract page https://arxiv.org/abs/2512.23382v2. A full rights notice is delivered at licenses/arxiv-2512.23382-CC-BY-4.0.txt. Characters: every retained excerpt is pure ASCII, so the delivered engine, XeLaTeX with its default Latin Modern text font, reports no missing character.
   \begin{lemma}\label{lem: the ineuqlaity in sec 3 (2)}
 For integers $r,k,\ell$ with $k\geq 3$ and $\ell\geq r\geq 3$, we have
     $$\frac{1}{2}\sum_{t=1}^{r-2}\binom{(k-1)\ell-1}{r-t-1}-\ell+\binom{\ell-1}{r-2}>0.$$
 \end{lemma}

 \begin{align}\label{3-2}
|E(\mh)|&\ge {k\ell'-1\choose r-1}(n-k\ell'+1)+{k\ell'-1\choose r}+\mathbb{I}_\ell\cdot\binom{k\ell'-1}{r-2}\\
&= \binom{k\ell'-1}{r-1}n+O_{r,k,\ell}(1)\nonumber.
\end{align}

 \begin{lemma}\label{lem: enough Berge common}
     There exists a constant $\alpha=\alpha(r,k,\ell)>0$ such that for every copy of Berge-$P_{\ell}$ in $\mathcal{H}$, say $\mathcal{P}$, there is $V_0\subseteq DV_\mh(\mathcal{P})$ with $|V_0|=
     \ell'$ and $|BCN(V_0)|\ge\alpha n$.
 \end{lemma}

 \begin{proof} Let $\mathcal{P}$ be  a copy of $\ber P_{\ell}$ in $\mh$. Set $D=DV_\mh(\mathcal{P})$ for short. Then $|D|=\ell+1$. For $0\le t\le r$, we set
         $$\mathcal{S}_t=\{E\in \mathcal{H}:|E\cap D|=t\}.$$
     Then $\mathcal{H}=\bigcup_{t=0}^{r}\mathcal{S}_t$ and $|\mathcal{S}_r|\leq \binom{\ell+1}{r}$.
     For $0\le t\le r-2$, we set $\mh_t=\{E\setminus D: E\in \mathcal{S}_t\}$ and $DE_\mh(\mathcal{P})|_{\mh_t}=\{E\setminus D: E\in DE_\mh(\mathcal{P})\cap \mathcal{S}_t\}$. Then $\mh_t$ is  $(r-t)$-uniform for $0\le t\le r-2$. Also $\mh_t':=\mh_t-DE_\mh(\mathcal{P})|_{\mh_t}$ is $\ber (k-1)P_{\ell}$-free for $0\le t\le r-2$; otherwise we can find a copy of $\ber kP_{\ell}$ in $\mh$.
     %Thus when $k=2$ we have
     %\begin{equation*}
      %   \begin{aligned}
       %      |E(\mh_t)|\leq & ex_{r-t}(n-\ell-1,\ber P_{\ell})+|DE_\mh(\mathcal{P})|_{\mh_t}|\\
        %     \leq & \binom{\ell}{r-t}\frac{n-\ell-1}{\ell}+\ell=\frac{1}{\ell}\binom{\ell}{r-t}n+O_{r,t,\ell}(1).
   %      \end{aligned}
    % \end{equation*}
    By the induction hypothesis, for $0\le t\le r-2$, we have
     \begin{equation}\label{lemm4-2}
         \begin{aligned}
             |\mh_t|\leq & \mathrm{ex}_{r-t}(n- \ell-1,\ber (k-1)P_{\ell})+|DE_\mh(\mathcal{P})|_{\mh_t}|\\
             \leq & \binom{(k-1)\ell'-1}{r-t-1}(n-(k-1)\ell'+1)+\binom{(k-1)\ell'-1}{r-t}+\binom{(k-1)\ell'-1}{r-t-2}+\ell\\
             \leq &  \binom{(k-1)\ell'-1}{r-t-1}n+O_{r,k,t,\ell}(1).
         \end{aligned}
     \end{equation}
    Note that    $|\mathcal{S}_0|=|\mh_0'|.$ For $1\le t\le r-2$ and $E\in \mh_t'$, define
          \begin{equation*}
         \begin{aligned}
         &C_t(E)=\{E'\in \mathcal{S}_t : E\subseteq E'\},~~c_t(E)=|C_t(E)|,\\
             &\mathcal{E}_{t,1}=\left\{E\in \mh_t' : c_t(E)\geq \binom{\ell'}{t}\right\},\\
             &\mathcal{E}_{t,2}=\left\{E\in \mh_t' : c_t(E)\leq \binom{\ell'}{t}-1\right\}.
         \end{aligned}
     \end{equation*}
     Then we have the following  claim.
     \begin{claim}\label{claim: et1 is small}
        For any $1\le t\le r-2$,  we may assume that $|\mathcal{E}_{t,1}|\leq \eta n$, where  $$ \eta=\frac{1}{2}\min_{t\in\{1,\ldots,r-2\}}\left\{\frac{\binom{(k-1)\ell'-1}{r-t-1}}{\binom{\ell+1}{t}-\binom{\ell'}{t}+1}\right\} >0.$$
     \end{claim}


    % \begin{claim}\label{claim: claim2}
         %Let $ \eta=\min_{t=1,2,\ldots,r-2}\{\frac{1}{2}\frac{\binom{(k-1)\ell'-1}{r-t-1}}{\binom{\ell+1}{t}-\binom{\ell'}{t}+1}\} >0$, and we may assume that $|\mathcal{E}_{t,1}|\leq \eta n$ for $t=1,2,\ldots,r-2$.
    % \end{claim}
    \noindent{\bf Proof of Claim \ref{claim: et1 is small}}
    Assume there is $1\le t\le r-2$ such that
 $|\mathcal{E}_{t,1}|>\eta n$. Let
  $E\in \mathcal{E}_{t,1}$ and set $C_t(E)=\{E_{1}',E_{2}',\ldots,E_{c_t(E)}'\}$.
   We construct a hypergraph $\mg_{t,E}=\{E_{j}'\setminus E : j=1,2,\ldots,c_t(E)\}$. Then $\mg_{t,E}$ is a $t$-graph. We use $N_{t,\ell}$ to denote the number of
   $t$-graphs with vertex set $D$. Since $|D|=\ell+1$, $|\mathcal{E}_{t,1}|>\eta n$ and $n$ is large enough, by the Pigeonhole Principle,  there exist $\mathcal{B}_{t,1}\subseteq \mathcal{E}_{t,1}$ with $|\mathcal{B}_{t,1}|\geq \frac{\eta n}{N_{t,\ell}}$ such that for any $E_1,E_2\in\mathcal{B}_{t,1}$ we have $\mg_{t,E_1}=\mg_{t,E_2}$. Note that $\mathcal{B}_{t,1}$ is also a $(r-t)$-hypergraph. Since  $\mh_t'$ is $\ber (k-1)P_{\ell}$-free, $\mathcal{B}_{t,1}$ is $\ber (k-1)P_{\ell}$-free. We have that $|V(\mathcal{B}_{t,1})|\ge
     \frac{\eta n}{2 \binom{(k-1)\ell'-1}{r-t-1}N_{t,\ell}}$; otherwise, by $\mathcal{B}_{t,1}$ being $\ber (k-1)P_{\ell}$-free and  the induction hypothesis, we have
     $$|\mathcal{B}_{t,1}|\leq \mathrm{ex}_{r-t}(|V(\mathcal{B}_{t,1})|,\ber (k-1)P_{\ell})\leq \frac{\eta n}{2 N_{t,\ell}}+O(1)<\frac{\eta n}{N_{t,\ell}},$$
      a contradiction.
      %Let $$U_{r-t}=\{v\in V(\mh) : \text{there exists $E\in\mathcal{B}$ such that $v\in E$}\}.$$
      When $t\geq 2$, for each $E\in \mathcal{B}_{t,1}$, we define
     $$S_{t,E}^{0}=\{E'\in \mathcal{S}_{t} : \text{$E\subseteq E'$ and there exist $u_1,u_2\in E'\cap D$ such that $d_{\mg_{t,E}}(u_i)=1,i=1,2$ } \},$$
     and set $S_{1,E}^{0}=\emptyset$. Then we have $|S_{t,E}^{0}|\leq \ell'$ for $t\geq 2$ by $|D|=\ell+1$ and $|S_{1,E}^{0}|=0$.
          For $3\leq t\leq r-2$ and $E\in \mathcal{E}_{t,1}$, we have
     $$|\mg_{t,E}|-|S_{t,E}^{0}|=c_t(E)-|S_{t,E}^{0}|\geq\binom{\ell'}{t}-\ell'\geq \binom{\ell'-1}{t},$$
     which implies $|V(\mg_{t,E}\setminus S_{t,E}^{0})|\ge \ell'-1$ by  $\ell'\ge r$ and $2\ell'\ge r+6$. Then there is $V_0'\subseteq V(\mg_{t,E}\setminus S_{t,E}^{0})$ with $|V_0'|= \ell'-1$ such that $V(\mathcal{B}_{t,1})\subseteq BCN(V_0')$. Since  $|V(\mg_{t,E})|\ge \ell'$ by $|\mg_{t,E}|=c_t(E)\geq \binom{\ell'}{t}$, there is $v_0\in V(\mg_{t,E})\setminus V_0'$ such that $V(\mathcal{B}_{t,1})\subseteq BCN(V_0)$, where $V_0=V_0'\cup\{v_0\}$. Then  $|BCN(V_0)|\ge |V(\mathcal{B}_{t,1})|\ge \frac{\eta n}{2 \binom{(k-1)\ell'-1}{r-t-1}N_{t,\ell}}$.



     Now we consider the case $t=2$. If $|S_{2,E}^{0}|=\ell'$, by the definition of $S_{2,E}^{0}$, $S_{2,E}^{0}$ is a match  in $D$. For each edge in $S_{2,E}^{0}$, we select one endpoint and thus we have  $V_0\subseteq V(S_{2,E}^{0})$ with $|V_0|= \ell'$ such that $V(\mathcal{B}_{2,1})\subseteq BCN(V_0)$. If $|S_{2,E}^{0}|\leq \ell'-1$, we have
     $$|\mg_{2,E}|-|S_{2,E}^{0}|\geq\binom{\ell'}{2}-(\ell'-1)\geq \binom{\ell'-1}{2}\ge \ell'$$ by $\ell'\ge 5$.
     So we have  $V_0\subseteq V(\mg_{2,E}\setminus S_{2,E}^{0})$ with $|V_0|= \ell'$ such that $V(\mathcal{B}_{2,1})\subseteq BCN(V_0)$.

     Finally, we consider the case $t=1$. Since $|\mg_{1,E}|\geq \binom{\ell'}{t}$,  $|V(\mg_{1,E})|\ge \ell'$. Then there is $V_0\subseteq V(\mg_{t,E})$ with $|V_0|= \ell'$ such that $V(\mathcal{B}_{1,1})\subseteq BCN(V_0)$.
     \par
     Above all, set $\alpha=\min_{t=1,\ldots,r-2}\{\frac{\eta}{2\binom{(k-1)\ell'-1}{r-t-1} N_{t,\ell}}\}>0$. If there is $1\le t\le r-2$ such that
 $|\mathcal{E}_{t,1}|>\eta n$, by the  discussion above, Lemma \ref{lem: enough Berge common} holds.\q

     %For the $\ber P_{\ell}$ $\mathcal{P}$ in $\mh$, if there exist $t\in\{1,2,\ldots,r-2\}$ such that $|\mathcal{E}_{t,1}|>\eta n$, we can find $\ell'$ vertices in $D$ and they have $|U_{r-t}|>\frac{\eta n}{2 \binom{(k-1)\ell'-1}{r-t-1}N_{t,\ell}}\geq \alpha n$ $\ber$common neighbors, otherwise Claim \ref{claim: claim2} holds. \q
Now we complete the proof of Lemma \ref{lem: enough Berge common}.
     By Claim \ref{claim: et1 is small} and (\ref{lemm4-2}), for $1\le t\le r-2$, we have
     \begin{equation}\label{lemma4-1}
     \begin{aligned}
     |\mathcal{S}_{t}|&=\sum_{E\in \mh_t}c_t(E)\leq |\mathcal{E}_{t,1}|\binom{\ell+1}{t}+|\mathcal{E}_{t,2}|\left(\binom{\ell'}{t}-1\right)\\
     & \leq \binom{\ell+1}{t}|\mathcal{E}_{t,1}|+(|\mh_t|-|\mathcal{E}_{t,1}|)\left(\binom{\ell'}{t}-1\right)\\
     &\leq  \binom{\ell+1}{t}|\mathcal{E}_{t,1}|+\left(\binom{(k-1)\ell'-1}{r-t-1}n-|\mathcal{E}_{t,1}|\right)\left(\binom{\ell'}{t}-1\right)+O_{r,k,t,\ell}(1)\\
     &=\left(\binom{\ell+1}{t}-\binom{\ell'}{t}+1\right)|\mathcal{E}_{t,1}|+\binom{(k-1)\ell'-1}{r-t-1}\left(\binom{\ell'}{t}-1\right) n+O_{r,k,t,\ell}(1)\\
     &\leq \left(\binom{\ell+1}{t}-\binom{\ell'}{t}+1\right)\eta n+\binom{(k-1)\ell'-1}{r-t-1}\left(\binom{\ell'}{t}-1\right) n+O_{r,k,t,\ell}(1)\\
     &\leq \binom{(k-1)\ell'-1}{r-t-1}\left(\binom{\ell'}{t}-\frac{1}{2}\right) n+O_{r,k,t,\ell}(1).
     \end{aligned}
     \end{equation}
     Thus when $k\geq 3$, by (\ref{3-2}) and (\ref{lemma4-1}), we have
     \begin{equation}\label{4-41}
         \begin{aligned}
             |\mathcal{S}_{r-1}|&=|\mh|-\sum_{t=0}^{r-2}|\mathcal{S}_t|-|\mathcal{S}_r|\\
             &\geq \binom{k\ell'-1}{r-1}n+O_{r,k,\ell}(1)-\left(\binom{(k-1)\ell'-1}{r-0-1}+\sum_{t=1}^{r-2}\binom{(k-1)\ell'-1}{r-t-1}\left(\binom{\ell'}{t}-\frac{1}{2}\right)\right)n\\
            & -O(r,k,t,\ell)(1)-\binom{\ell+1}{r}\\
            &\geq\left(\binom{k\ell'-1}{r-1}-\sum_{t=0}^{r-2}\binom{(k-1)\ell'-1}{r-t-1}\binom{\ell'}{t}\right)n+\alpha_0 n-O_{r,k,t,\ell}(1),
         \end{aligned}
     \end{equation}
      where $
     \alpha_0=\frac{1}{2}\sum_{t=1}^{r-2}\binom{(k-1)\ell'-1}{r-t-1}>0$.
          Fix a vertex $v\in V(\mh)\setminus D$, we can construct an $(r-1)$-uniform hypergraph $\mathcal{D}_v=\{E\setminus \{v\}: E\in \mathcal{S}_{r-1} \text{~and $v\in E$}\}$. %For $u\in D$, recall that $E_{\mathcal{D}_v}(u)$ denote the collection of hyperedges of $\mathcal{D}_v$ that contains $u$.
    We set $$S_{r-1,v}^0=\{E\in \mathcal{S}_{r-1}: \text{~$v\in E$ and there exist $u_1,u_2\in E\cap D$ with $d_{\mathcal{D}_v}(u_i)=1$, $i=1,2$}\},$$
    and $$\mathcal{S}^+_{r-1,v}=\{E\in \mathcal{S}_{r-1}\setminus\mathcal{S}_{r-1,v}^0:\text{$v\in E$}\}.$$
    For every two vertices $u_1,u_2\in \left(\bigcup_{E\in S^+_{r-1,v}}E\right)\cap D$, there exist two different hyperedges $E_1,E_2\in E(\mathcal{D}_{v})$ with $\{u_i,v\}\subseteq E_i$ for $i=1,2$, which implies $v\in BCN(\{u_1,u_2\})$.

   Denote $\mathcal{S}_{r-1}^0=\bigcup_{v\in V(\mh)\setminus D} S_{r-1,v}^0$ and $\mathcal{S}^+_{r-1}=\bigcup_{v\in V(\mh)\setminus D}\mathcal{S}_{r-1,v}^+$. Then $S^+_{r-1}=\mathcal{S}_{r-1}\setminus S_{r-1}^0$. Note that for every $v\in V(\mh)\setminus D$,  $|S_{r-1,v}^0|\leq \ell'$. Then $|\mathcal{S}_{r-1}^0|\leq \ell' n$ and $|\mathcal{S}_{r-1}^+|\geq |\mathcal{S}_{r-1}|-\ell' n$.

     For a vertex $v\in V(\mh)\setminus D$, we set $$N_{\mathcal{S}_{r-1}^+,D}(v)=\{u\in D: \text{there exist $E\in \mathcal{S}_{r-1}^+$ such that $\{u,v\}\subseteq E$}\},$$ and  $d_{\mathcal{S}_{r-1}^+,D}(v)=|N_{\mathcal{S}_{r-1}^+,D}(v)|$.
    Let $B=\{v\in V(\mh)\setminus D:d_{\mathcal{S}_{r-1}^+,D}(v)\geq \ell'\}$. Then we have
     \begin{equation*}
         \begin{aligned}
             \left|\mathcal{S}_{r-1}^+\right|&\leq \binom{\ell+1}{r-1}|B|+\binom{\ell'-1}{r-1}(n-|B|)\\
             &= \binom{\ell'-1}{r-1}n+\left( \binom{\ell+1}{r-1}-\binom{\ell'-1}{r-1} \right)|B|.
         \end{aligned}
     \end{equation*}
     Note that  $\ell'\geq r\geq 3$ and $k\geq 3$. By (\ref{4-41}) and $|\mathcal{S}_{r-1}^+|\geq |\mathcal{S}_{r-1}|-\ell' n$, we have
      \begin{equation*}
         \begin{aligned}
             &\left( \binom{\ell+1}{r-1}-\binom{\ell'-1}{r-1} \right)|B|\geq \left| \mathcal{S}_{r-1}^+\right|-\binom{\ell'-1}{r-1}n\\
             &\geq \left(\binom{k\ell'-1}{r-1}-\sum_{t=0}^{r-2}\binom{(k-1)\ell'-1}{r-t-1}\binom{\ell'}{t}-\ell'-\binom{\ell'-1}{r-1} \right)n
             +\alpha_0 n-O_{r,k,\ell}(1)   \\
             &= \left(\binom{\ell'-1}{r-2}-\ell'+\alpha_0 \right)n-O_{r,k,\ell}(1).
         \end{aligned}
     \end{equation*}
     By Lemma \ref{lem: the ineuqlaity in sec 3 (2)}, there exists $\beta=\beta(r,k,\ell)>0$ such that $|B|\geq \beta n$.
     Since there are $\binom{\ell+1}{\ell'}$ sets with size $\ell'$ in $D$, there are $B_0\subseteq B$ with  $|B_0|\ge\frac{\beta }{\binom{\ell+1}{\ell'}}n$ and $V_0\subseteq D$ with $|V_0|=\ell'$ such that for every $v\in B_0$ and $u\in V_0$, there exists hyperedges $E\in \mathcal{S}_{r-1}^+$ with $\{u,v\}\subseteq E$.
So for every $v\in B_0$ and $u\in V_0$, $u\in \left(\bigcup_{E\in \mathcal{S}_{r-1,v}^+}E\right)\cap V_0$.
     Then for every two vertices $u_1,u_2\in V_0$ and every vertex $v\in B_0$, we have $u_1,u_2\in  \left(\bigcup_{E\in \mathcal{S}_{r-1,v}^+}E\right)\cap V_0$. Thus $B_0\subseteq BCN(V_0)$.
    By setting $\alpha=\frac{\beta}{\binom{\ell+1}{\ell'}}$, we finished the proof.
 \end{proof}

\end{document}
